\section{四问递进关系与统一接口}

四个问题共享“任务--区域--时间--功率”口径，但每一步只接收上一步已经核验的输出。问题一先完成需求预测并对真实到达任务作可行调度；问题二在不改变任务时序语义的前提下，引入区域电价、碳强度和可用新能源，得到碳感知任务调度；问题三固定问题二的任务负荷，优化储能充放电和区域购售电；问题四再将任务包络、储能状态和电网边界放入同一顺序耦合框架，并用独立峰值扫描刻画削峰代价。预测值只用于预测评价和压力分析，不能替代问题一收尾窗口的实际到达任务。

任务 $i$ 的 GPU 需求、持续时间和到达时刻分别记为 $g_i$、$d_i$ 和 $a_i$，区域和任务类型为 $r_i$、$k_i$。统一工作量为
\begin{equation}
  W_{rkt}=\sum_{i:r_i=r,\,k_i=k,\,a_i=t}g_i\frac{d_i}{60},
  \qquad
  L_{rt}=L^{\mathrm{nonAI}}_{rt}+\mathrm{PUE}_{r}P^{\mathrm{AI}}_{rt}.
\end{equation}
其中 $L_{rt}$ 是储能模型的设施负荷输入。每一问均输出与基线同类的任务记录或功率记录，并以独立审计文件中的字段为主张来源。

\section{问题二：碳感知全时域任务调度}
\label{sec:q2-proposal}

\subsection{模型选择与变量}

问题二的核心是带释放时间、不可抢占和跨区域时延的在线组合调度。若直接建立 0--2399 小时的单体 MILP，变量规模会随 50,000 个任务和分钟级开始时刻急剧增长；纯 FIFO 又无法利用碳强度和新能源的时间差。因此采用确定性全时域滚动影子交换启发式，在每个滚动窗口内先构造可行候选，再比较有限交换。设 $x_{ir}\in\{0,1\}$ 表示任务区域，$s_i$ 和 $e_i=s_i+d_i$ 为开始和完成时刻，$y_{irt}$ 表示任务在时刻 $t$ 占用区域 $r$。核心可行域写为
\begin{equation}
  \sum_r x_{ir}=1,\quad s_i\ge a_i,\quad
  y_{irt}=1\Longleftrightarrow x_{ir}=1,\ s_i\le t<e_i,
\end{equation}
\begin{equation}
  \sum_{i,k}g_i y_{irt}\le G_r,\quad
  P^{\mathrm{AI}}_{rt}\le \overline P^{\mathrm{IT}}_r,
  \quad \mathrm{PUE}_rP^{\mathrm{AI}}_{rt}+P^{\mathrm{nonAI}}_{rt}\le \overline P^{\mathrm{fac}}_r,
\end{equation}
并对每个任务施加静态网络时延和 SLA 上限。对可行候选 $x$，窗口评分接口为
\begin{equation}
  J_t(x)=\omega_C\widehat C_t(x)+\omega_E\widehat E_t(x)
  +\omega_L\widehat D_t(x)-\omega_R\widehat U_t(x)
  +\sum_r\lambda_{r,t}\,[P_{r,t}(x)-\overline P_r]_+,
\end{equation}
其中带帽量按同一窗口的基线归一化，$\lambda_{r,t}$ 是容量或功率压力的影子项。算法沿固定任务顺序进行一次有界交换，只有在 $J_t$ 改善且硬约束复核通过时才接受交换；分钟级精确扫描仅在小时网格无法解析柔性任务的开始位置时触发。

\subsection{基线、结果与稳健性}

FIFO--local-first 基线使用完全相同的 50,000 个任务、约束、功率映射、静态时延矩阵和输出字段。候选和基线均通过任务唯一性、不可抢占、释放时间、容量、功率、时延与收尾边界审计。候选覆盖率为 $1.000000$，SLA 违约率为 $0.000000$。相对 FIFO，设施运行成本变化为 $-1.0149373918772464\%$，碳排变化为 $-1.929297238773453\%$，加权平均网络时延变化为 $+4.18976\,\mathrm{ms}$。这三个数字应同时呈现，不能把成本和碳排改善外推为所有指标均改善；候选平均时延为 $9.25822\,\mathrm{ms}$，基线为 $5.06846\,\mathrm{ms}$。

固定调度压力探针共 8 组，块级成本变化范围为 $[-1.0947051309\%,-0.9104543589\%]$，碳排变化范围为 $[-2.0231818682\%,-1.8692647719\%]$；确定性复跑通过。上述压力探针只度量固定计划的余量暴露，不等同于重新优化后的性能保证。由于主算法是有界一遍交换启发式，本文不使用“全局最优”“MILP 最优”或“唯一最优策略”等表述；本问也不优化储能和电网交易，这些变量留给问题三、四。

\section{问题三：多区域滚动储能协同}
\label{sec:q3-proposal}

\subsection{二进制储能 MILP}

问题三接收问题二的任务调度和逐时设施负荷，不再重排任务。区域 $r$ 在滚动块 $b$ 的时刻 $t$ 设充电功率 $p^{\mathrm{ch}}_{rt}$、放电功率 $p^{\mathrm{dis}}_{rt}$、荷电状态 $q_{rt}$、购电 $p^{\mathrm{buy}}_{rt}$、售电 $p^{\mathrm{sell}}_{rt}$、弃新能源 $p^{\mathrm{curt}}_{rt}$，并设 $u_{rt}\in\{0,1\}$ 表示充放电模式。能量和互斥接口为
\begin{equation}
  q_{r,t+1}=q_{rt}+\eta_c p^{\mathrm{ch}}_{rt}
  -p^{\mathrm{dis}}_{rt}/\eta_d,
  \quad 0\le q_{rt}\le E_r,
\end{equation}
\begin{equation}
  0\le p^{\mathrm{ch}}_{rt}\le u_{rt}\overline P^{\mathrm{ch}}_r,
  \quad 0\le p^{\mathrm{dis}}_{rt}\le (1-u_{rt})\overline P^{\mathrm{dis}}_r.
\end{equation}
区域购售电边界、终端 SOC 下界和电量平衡共同保证物理可行性。以成本、购电碳排、区域净购电峰值、负荷波动和弃新能源的归一化加权和为目标，得到每个滚动块的二进制 MILP；无储能新能源优先平衡作为同口径基线。

每个区域使用 14 个 168 小时块和 1 个 55 小时收尾块，六个区域共 90 次成功求解，精确覆盖 $0$--$2406$ 小时。块间 SOC 按下界和 carry-forward 传递，不能写成每个块均回到初始 SOC；该滚动证据也不等价于一个 2407 小时单体 MILP。

\subsection{结果、压力探针与审计边界}

六区域滚动汇总相对于无储能基线的成本差为 $-38,899,969.87293261\,\mathrm{CNY}$，碳排差为 $-8.8505625719\,\mathrm{tCO_2}$；负成本表示声明的运行成本公式下存在净售电收益。四类由观测数据派生的 72 小时情景压力探针中，候选与基线硬约束审计 48/48 通过。设施负荷由非 AI 负荷、AI 功率映射和区域 PUE 独立复算，最大设施残差为 $5.044995509706496\times10^{-5}\,\mathrm{MW}$，因此只能称“在舍入残差范围内复现”，不能称精确恒等。

完整审计共 270 项，其中 269 项通过。永久排除的全周期 LP 可扩展性探针在 RegionF 出现 $116.60975349482445\,\mathrm{MW}$ 同时充放电，故其任何指标不进入论文主张。滚动二进制 MILP 的压力结果支持可行性和趋势判断，但不支持全时域全局最优或 fallback 性能主张；本次 fallback 使用块数为 0。

\section{问题四：顺序耦合的算--储--电权衡}
\label{sec:q4-proposal}

\subsection{固定包络下的联合接口}

问题四首先分别由问题二候选和 FIFO 生成设施负荷包络，再在绝对小时 $2328$--$2399$ 上对六个区域构造 72 小时二进制储能 MILP。任务记录、储能状态和电网购售电不在同一模型中重新联合搜索，因此保留了“任务调度先决、储能随后”的顺序接口。对购电 $b_{rt}$、售电 $v_{rt}$ 设模式变量 $h_{rt}\in\{0,1\}$，例如
\begin{equation}
  0\le b_{rt}\le M h_{rt},\qquad
  0\le v_{rt}\le M(1-h_{rt}),
\end{equation}
并将问题三的 SOC、充放电互斥、区域购售边界和终端下界一并保留。系统峰值用上界变量 $\pi$ 连接：
\begin{equation}
  \pi\ge\sum_r b_{rt},\quad t=1,\ldots,72.
\end{equation}
候选目标为成本、碳排、峰值、弃电和微小吞吐正则项的离散复合目标；在观测、峰段电价、高碳、低新能源和联合压力五类情景下分别求解，并对硬约束独立复算。

\subsection{低新能源结果与独立峰值扫描}

低新能源 72 小时情景中，候选任务完成率为 $1.000000$、SLA 违约率为 $0.000000$，相对无储能 FIFO 基线的成本差为 $-1,683,494.3771166522\,\mathrm{CNY}$，碳排差为 $-189.1085063591487\,\mathrm{tCO_2}$，新能源利用率提高 $0.0115804947$，平均时延为 $9.0329740013\,\mathrm{ms}$。联合压力情景的成本差和碳排差分别为 $-1,683,494.3771166522\,\mathrm{CNY}$ 与 $-225.6074388960753\,\mathrm{tCO_2}$；这些差值同时包含上游任务包络和储能调度效应，不能拆成纯储能因果效应。

为识别峰值权衡，另在低新能源 24 小时包络上扫描 17 个离散峰值权重。零权重系统峰值为 $291.60330044\,\mathrm{MW}$，权重 $0.0003$ 的选定点峰值为 $197.08826651333345\,\mathrm{MW}$，削峰 $94.51503392666643\,\mathrm{MW}$；离散复合目标边际估计为 $2.35298451326919\times10^{-8}$ composite-objective/MW。该边际量只用于相邻离散 MILP 点的比较，不是 LP 对偶变量或连续影子价格。

\subsection{四问整体评价与适用边界}

四问形成可复现的“预测--可行调度--滚动储能--顺序耦合”链路，主模型和基线始终使用同类输出，结果由冻结的 JSON/CSV 和独立审计定位。局限性也随链路保留：问题二的静态时延未建模拥塞、迁移能耗和迁移费用；问题三按区域滚动且不含区域间潮流；问题四固定任务包络、只覆盖 72 小时，并将峰值探针与主情景分开。上述边界决定了本文结论是给定数据和约束下的可行策略比较与机制证据，而非对全周期联合最优的证明。
